% !PRINTED 711
\weberpage{711}
%<<CONT
is, and accordingly the functions $f$ are also to be distinguished
into $f_1$ and $f_2$, so that

\weq{11}{f_1(r^{-s}) = f_1(r^s), \quad f_2(r^{-s}) = -\,f_2(r^s)}
holds. Likewise we distinguish the expressions (8), (9) as $X_1$
and $X_2$, and if both expressions are added, there results, by
(11):
\weq{12}{\begin{aligned}
X_1 &= -\frac{1}{2m}\sum_{1,\,m-1}^{s} f_1(r^s)\,\log\left(\sin\frac{s\pi}{m}\right)^{2}, \\
X_2 &= -\frac{i\pi}{m^2}\sum_{1,\,m-1}^{s} s\,f_2(r^s).
\end{aligned}}

\begin{center}
\S.\,195.\\[4pt]
\gesperrt{On the class number in the real field contained in
$\Omega_m$.}
\end{center}

Our general principles can also be applied to determine the class
numbers in the divisors of the field $\Omega_m$.

In particular the formulae (6), (7), \S.\,192 are applicable if we
know the degrees of the prime ideals in such a divisor. The question
of the ratio of the class numbers in the divisors of a field
$\Omega_m$ to the class number in $\Omega_m$ itself is of great
interest, and it shall be treated here for the case that the field
concerned is the real field $H_m$ contained in $\Omega_m$, which we
have indeed already investigated in \S.\,176 f.\footnote{Cf.
\gesperrt{Kummer}, \glqq Bestimmung der Anzahl u.\,s.\,f.\grqq{}.
\gesperrt{Crelle's Journ.}, Bd.\,40 (1849).} It was found there that
$q$ in the field $H_m$ is the $\tfrac12\mu$th power of a prime of
the first degree, that the other primes $p$ split into $e_1$ prime
factors of degree $f_1$, so that

\weq{1}{f_1 e_1 = \tfrac{1}{2}\mu,}
and that two kinds of primes $p_1$, $p_2$ are to be distinguished,
so that for the primes of the first kind $f_1 = \tfrac12 f$,
$e_1 = e$, and for those of the second kind $f_1 = f$,
$e_1 = \tfrac12 e$.

For the field $H_m$ we accordingly obtain from \S.\,192, (6), (7)
the class number $h_1$ by the formula:
%CONT>>
